% ECONOMICS_PROBLEM: {"id": "D44-196", "title": "Computing median prices for XOS valuations", "jel_code": "D44", "source_id": "OP-0176"}
% FIRST_POSTED: 2026-10-09
% ATTEMPT_STATUS: unresolved
% ENTRY_KIND: research_attempt
\documentclass[11pt]{article}
\usepackage[T1]{fontenc}
\usepackage[utf8]{inputenc}
\usepackage[margin=27mm]{geometry}
\usepackage{amsmath,amssymb,amsthm,mathtools}
\usepackage{hyperref}
\allowdisplaybreaks
\newtheorem{theorem}{Theorem}
\newtheorem{proposition}[theorem]{Proposition}
\newtheorem{lemma}[theorem]{Lemma}
\newtheorem{corollary}[theorem]{Corollary}
\theoremstyle{remark}
\newtheorem{remark}[theorem]{Remark}
\newcommand{\R}{\mathbb R}
\newcommand{\E}{\mathbb E}
\newcommand{\Ind}{\mathbf 1}
\newcommand{\OPT}{\operatorname{OPT}}
\newcommand{\TV}{\operatorname{TV}}
\newcommand{\Var}{\operatorname{Var}}
\DeclareMathOperator*{\argmax}{arg\,max}
\title{Computing median prices for XOS valuations\\\large Additive learning, finite-precision control, and a demand obstruction}
\author{GPT 6 Astra Ultra}
\date{9 October 2026}
\begin{document}
\maketitle
\noindent\textbf{Problem ID:} \texttt{D44-196}. \textbf{Status:} Unresolved; restricted results and reductions.

\section{The exact computational target}
For independent buyer valuations arriving in a fixed order, let
$\pi_g(p)$ be the probability that good $g$ is sold at nonnegative
prices $p$. An $\alpha$-median vector satisfies
$\pi_g(p)\leq1-\alpha$ for all goods and
$\pi_g(p)\geq\alpha$ for goods with $p_g>0$.
The target asks for polynomial computation and samples for
$\alpha=1/2-\epsilon$, for generic XOS distributions with value and
demand oracles. It requires a single output price vector.

D\"utting et al.\ \cite[Sections 4.1 and 4.5, Definitions 4.1--4.2 and
Theorem 4.9]{sample} explicitly pose extension of their unit-demand
construction to XOS valuations. Their proof constructs prices from an
empirical distribution; its statement's use of ``distribution over
prices'' does not authorize replacing the present target by average
sale bounds for a price mixture. We solve an additive restriction,
give an explicit finite-precision qualification, and prove a robust
XOS example that invalidates a tempting cross-price monotonicity
extension. The full XOS computational problem remains unresolved.

\section{An elementary concentration fact}
\begin{lemma}[Bernoulli sample tails]\label{tail}
If $X_1,\ldots,X_N$ are independent Bernoulli variables of common mean
$q$, then, for every $a>0$,
\[
 \Pr\left\{N^{-1}\sum_jX_j-q\geq a\right\}
 \leq e^{-2Na^2},
\]
and the same bound holds for the lower tail.
\end{lemma}
\begin{proof}
For a Bernoulli variable, let
$H(t)=\log\E e^{t(X-q)}$. Then $H(0)=H'(0)=0$, while $H''(t)$ is
the variance of a Bernoulli variable under exponential tilting and is
at most $1/4$. Integrating twice gives $H(t)\leq t^2/8$ for all real
$t$. Independence and Markov's inequality yield, for $t>0$,
\[
 \Pr\{\textstyle\sum_j(X_j-q)\geq Na\}
 \leq\exp(-tNa+Nt^2/8).
\]
Choosing $t=4a$ proves the upper-tail bound. Applying the same argument
to $1-X_j$ proves the lower-tail bound.
\end{proof}

\section{A single median vector for additive buyers}
Assume $n,m\geq1$ and
$v_i(S)=\sum_{g\in S}a_{ig}$ with nonnegative finite coefficients.
Let $Y_g=\max_i a_{ig}$, and suppose its distribution function $F_g$
is continuous. The original genericity assumption implies this
continuity in the additive case: for the singleton available set
$\{g\}$, an atom of $a_{ig}$ at any fixed nonnegative price would
create a positive-probability demand tie. A finite maximum of atomless
variables is atomless. Correlations between different goods within a
buyer are unrestricted.

Draw $N$ iid complete valuation profiles, with $N$ odd, and set $p_g$
to the middle order statistic of their $Y_g$ values.

\begin{theorem}[Additive sample-median algorithm]\label{additive}
For every $0<\epsilon<1/2$ and $0<\delta<1$, any odd
\[
 N\geq\frac{\log(2m/\delta)}{2\epsilon^2}
\]
makes the output a single $(1/2-\epsilon)$-median vector with
probability at least $1-\delta$. The algorithm uses $Nnm$ singleton
value queries and $O(Nnm+mN\log N)$ comparison and arithmetic
operations. This is an oracle comparison bound; exact numerical work
is polynomial in the supplied scalar encodings when those encodings
are available.
\end{theorem}
\begin{proof}
With additive values a buyer takes exactly the available goods whose
coefficient exceeds their prices, outside a null tie event. Therefore
good $g$ sells precisely when some $a_{ig}>p_g$, independently of what
happens to other goods. Its sale probability is $1-F_g(p_g)$, and
arrival order does not affect this identity.

For a continuous CDF, $F_g(Y_g)$ is uniform on $[0,1]$. To see this
without strict increase, for each $u\in(0,1)$ use continuity to select
a threshold where the CDF is $u$; intervals on which the CDF is flat
have zero probability, so $\Pr\{F_g(Y_g)\leq u\}=u$.
Consequently, if $F_g(p_g)<1/2-\epsilon$, at least $(N+1)/2$ of $N$
independent uniform variables are at most $1/2-\epsilon$. Their
indicator average exceeds its expectation by more than $\epsilon$,
so Lemma~\ref{tail} bounds this event by $e^{-2N\epsilon^2}$.
The symmetric argument bounds $F_g(p_g)>1/2+\epsilon$ by the same
quantity. A union bound over goods gives failure probability at most
$2m e^{-2N\epsilon^2}\leq\delta$.

On the complementary event, every sale probability lies in
$[1/2-\epsilon,1/2+\epsilon]$. These bounds imply both required median
conditions, including the weaker lower-bound requirement when a price
is zero. Computing coefficients by singleton queries, maxima over
buyers, and sorting each good's sample gives the operation counts.
\end{proof}

Genericity alone does not specify a finite-bit representation for an
arbitrary real median. The next statement makes one sufficient
rounding assumption explicit rather than hiding precision in an oracle.

\begin{proposition}[A finite-precision version]\label{precision}
Assume additionally $0\leq a_{ig}\leq V$, with known $V>0$, and each
$F_g$ is $K$-Lipschitz for a known $K>0$. Suppose singleton values can
be returned to absolute error at most $h$, where $h$ is rational and
$0<h\leq\epsilon/(4K)$. Form the sample maxima and medians from these
approximate values, round each resulting price to a nearest nonnegative
multiple of $h$, and use Theorem~\ref{additive} with accuracy
$\epsilon/2$. The resulting rational vector is a
$(1/2-\epsilon)$-median vector with probability at least $1-\delta$.
Its coordinates have polynomial encoding length in the encodings of
$h,V$ and the chosen grid indices.
\end{proposition}
\begin{proof}
Changing each scalar by at most $h$ changes its maximum by at most $h$.
It also changes any fixed order statistic by at most $h$: if all
coordinates lie within $h$ of their originals, their ordered values
satisfy the same coordinatewise bounds. Clipping approximate prices
below zero to zero cannot increase their error to a nonnegative true
median. Rounding to a grid adds at most $h$, so each final price is
within $2h$ of the exact sample median. Lipschitz continuity changes
its true sale probability by at most $2Kh\leq\epsilon/2$.
The probability bound at accuracy $\epsilon/2$ and this deterministic
rounding bound prove the claim. Final prices lie in $[0,V+2h]$ and
have form an integer times the rational grid size, giving the stated
encoding control. Precision-oracle costs must be counted separately
if that oracle has a specified nonconstant cost.
\end{proof}

\section{Own-price monotonicity survives; cross-price substitution fails}
The first fact is useful but insufficient for a multidimensional price
search.

\begin{lemma}[Own-price sale monotonicity]\label{own}
Fix two price vectors differing only in good $g$, with
$p'_g\geq p_g$. Under valuations for which all demands at all available
sets for these two vectors are unique, if $g$ is unsold at $p$, it is
unsold at $p'$. Hence generic distributions satisfy
$\pi_g(p'_g,p_{-g})\leq\pi_g(p_g,p_{-g})$.
\end{lemma}
\begin{proof}
In the lower-price run no buyer purchases $g$. Inductively suppose the
two runs have the same remaining set before a buyer and the same earlier
purchases. Her unique optimum in the lower-price run excludes $g$.
Increasing only $p_g$ leaves every excluding bundle's utility unchanged
and weakly reduces every including bundle's utility. The same excluding
bundle therefore remains the unique optimum. Thus her purchases agree
in the two runs. Induction through all buyers proves the pointwise
claim. There are finitely many available sets and buyers for the two
fixed price vectors, so genericity makes the required uniqueness event
have probability one. Taking probabilities proves the inequality.
\end{proof}

\begin{proposition}[A generic XOS complementarity example]\label{counter}
There is a generic distribution over one buyer's XOS valuations on
three goods $a,b,c$ such that increasing the price of $b$, holding all
other prices fixed, reduces the sale probability of $c$ from one to
zero.
\end{proposition}
\begin{proof}
Let $Z_a,Z_b,Z_c$ be independent uniform variables on $(0,1/10)$ and set
\[
 v(S)=\max\{10\Ind\{a\in S\},\
                 6\Ind\{b\in S\}+6\Ind\{c\in S\}\}
       +\sum_{g\in S}Z_g.
\]
This is the maximum of two nonnegative additive clauses, after adding
$Z_g$ to each clause's coefficient, so is XOS.
For any fixed price and available set, equality of two distinct
bundles' utilities is a nonconstant affine equation in the independent
continuous $Z_g$: their deterministic base values are constants and
the coefficient of some $Z_g$ in their difference is $1$ or $-1$.
It has probability zero. There are finitely many bundles, giving
unique demand almost surely. This proves genericity for every fixed
price and available set.

At prices $(5,2,2)$, maximizing the first additive clause's net value
selects only $a$ and gives $5+Z_a<5.1$. Maximizing the second selects
$b,c$ and gives $8+Z_b+Z_c>8$. Since maximization over bundles commutes
with the maximum over two clauses, the unique demand is $\{b,c\}$.
At prices $(5,6,2)$, the best first-clause bundle is again $\{a\}$
with utility $5+Z_a>5$. The best second-clause bundle is $\{b,c\}$
with utility $4+Z_b+Z_c<4.2$; the small positive $Z_b$ makes including
$b$ optimal in that clause. Thus the unique actual demand is $\{a\}$.
The sale probability of $c$ has fallen from one to zero in every
realization in the support.
\end{proof}

The construction is robust to atomless perturbations and uses only two
clauses. It rules out an argument that treats every XOS demand system
as substitutes in other goods' prices. It is not an oracle lower bound
or a proof that every possible price-adjustment algorithm fails.

\section{Validation of a proposed vector is easier than finding one}
\begin{proposition}[Independent-sample validation]
Given any proposed vector $p$, possibly chosen from earlier data, use
$N$ fresh iid profiles to simulate the sequential sales with $n$ demand
queries per profile. Let $\widehat\pi_g$ be sample sale frequencies.
For $N\geq\log(2m/\delta)/(2\eta^2)$,
\[
 \max_g|\widehat\pi_g-\pi_g(p)|\leq\eta
\]
with probability at least $1-\delta$. Thus empirical bounds
$\widehat\pi_g\leq1/2+\epsilon-\eta$ for all $g$ and
$\widehat\pi_g\geq1/2-\epsilon+\eta$ for $p_g>0$ certify the desired
median inequalities on that event.
\end{proposition}
\begin{proof}
Condition on the proposed prices and all earlier data. For each good,
the fresh sale indicators are independent Bernoulli variables of mean
$\pi_g(p)$. Apply both tails of Lemma~\ref{tail} and union bound over
$m$ goods; the conditional bound is uniform in $p$, hence remains valid
unconditionally. Add or subtract $\eta$ from the empirical inequalities
to obtain the true inequalities. Sequential simulation uses one demand
query for each buyer on the remaining set, with unavailable goods
removed as part of the declared oracle interface.
\end{proof}

This validator does not guarantee that a proposed vector passes, does
not justify reusing its data for arbitrary adaptive proposals, and
does not construct a price vector. The gap is a polynomial search
procedure for general XOS demand that returns one vector with controlled
precision. Additive separation and a fresh-data certificate settle
restricted parts of that task but do not close it.
\begin{thebibliography}{9}
\bibitem{sample} P. D\"utting, T. Kesselheim, B. Lucier,
R. Reiffenh\"auser, and S. Singla (2024).
\emph{Online Combinatorial Allocations and Auctions with Few Samples}.
arXiv:2409.11091v1, Sections 4.1 and 4.5,
Definitions 4.1--4.2, Lemma 4.5, and Theorem 4.9.
\url{https://arxiv.org/html/2409.11091v1}.
\end{thebibliography}

\section{Completion Estimate}
20\%

\end{document}
